% ch05_u04 —— 原书 p.220--p.229（5.4.8 -- 5.6.11）
% 本单元从 5.4.8 起，末尾停在 5.6.11 的中途；尾部接 p.230。

\SourcePage{220}

设 $\mathit{Vac}'\in\mathcal{M}(\mathfrak{g}\otimes K,G(O))'$ 为扭真空模。

\bdp{5.4.8} \emph{定理.} 对任意对象 $P\in\mathcal{P}$，对象
$P\circledast\mathit{Vac}'\in D(\mathfrak{g}\otimes K,G(O))'$ 同构于若干份
$\mathit{Vac}'$ 的直和。\footnote{特别地，它是单个 Harish-Chandra 模，而非这样的复形。}

这一定理等价于 5.2.14。事实上，由 7.8.5 的（335）与 7.14.1，有
$(\tilde{\mathfrak{g}}\otimes K,G(O))$-模的典范等同
\begin{equation}\tag{282}
H^{i}(P\circledast\mathit{Vac}')=H^{i}(\mathcal{GR},P\lambda_{\mathcal{L}}^{-1}).
\end{equation}

这里 $P\lambda_{\mathcal{L}}^{-1}:=P\otimes\lambda_{\mathcal{L}}^{-1}$。有兴趣的读者可以
直接去看这一定理的证明，它可在 ??? 中找到。
% [Tutor Note: 原书 p.220 此处交叉引用印作「???」，底本未补，译文照录。]

\bdp{5.4.9} 我们需要把 $\operatorname{Aut}O$ 的对称性纳入上述设定。回忆（见 4.6.6）
$\operatorname{Aut}O$ 在 $G(K)$ 上的作用提升为 $\operatorname{Aut}_ZO$ 在
$\widetilde{G(K)}$ 上的作用，后者保持 $G(O)$。所以我们处在 7.9.5 的设定中。
\footnote{我们的 $(\operatorname{Der}O,\operatorname{Aut}^{0}_{Z}O)$ 与
$(\widetilde{G(K)},G(O))$ 就是 7.9.5 的 $(\mathfrak{l},P)$ 与 $(G',K)$。}

设 $D_{HC}$ 为 7.9.5 中定义的 Harish-Chandra 复形的导出范畴。这是一个
t-范畴，其心 $\mathcal{M}_{HC}$ 等于对
$(\operatorname{Der}O\ltimes\widetilde{\mathfrak{g}\otimes K},
\operatorname{Aut}^{0}_{Z}O\ltimes G(O))$ 的 Harish-Chandra 模的范畴
（我们假设中心 $\mathfrak{C}\subset\tilde{\mathfrak{g}}\otimes K$ 按标准方式作用）。
% [Tutor Note: 原书 p.220 此处「中心」印作 fraktur C（$\mathfrak{C}$），按上下文
% 应为 $\mathfrak{Z}$；底本如此，译文照录。]

对 $(G(K),G(O))$ 的 $(\operatorname{Der}O,\operatorname{Aut}^{0}_{Z}O)$-等变
Hecke 范畴（见 7.9.2）以全幺半子范畴的形式包含相应的
$(\operatorname{Der}O,\operatorname{Aut}^{0}O)$-等变范畴 $\mathcal{H}$ 与
$\mathcal{H}^{c}$。于是由 7.9.5，$D_{HC}$ 是 $\mathcal{H}$-模，从而是
$\mathcal{P}$-模。

我们需要对设定稍作改变。按惯例，设 $\mathfrak{Z}$ 为 $\mathfrak{g}\otimes K$ 的
完备化扭泛包络代数的中心，$\mathfrak{z}$ 为扭真空模 $\mathit{Vac}'$ 的自同态环；
有显然的代数态射 $e:\mathfrak{Z}\to\mathfrak{z}$。设 $D_{HC\mathfrak{z}}$ 为 7.9.8 中
定义的相应 Harish-Chandra 复形的导出范畴（另见
\SourcePage{221}
7.9.7(iii)）\footnote{我们的 $D_{HC\mathfrak{z}}$ 就是 7.9.8 的 $D^{e}_{HC\,A}$。
在 7.9.8 中，$\mathfrak{Z}$ 表示中心里 $G(K)$-不变的元素之集，但由 3.7.7(ii)，
中心的所有元素都是 $G(K)$-不变的。}。这是一个 t-范畴，其心
$\mathcal{M}_{HC\mathfrak{z}}$ 等于被 $\operatorname{Ker}e$ 消灭的
Harish-Chandra 模的范畴。

设 $\mathcal{H}_{\mathfrak{z}}$ 是 $(G(K),G(O))$ 的
$(\operatorname{Der}O,\operatorname{Aut}^{0}_{Z}O)$-等变 Hecke 范畴的
$\mathfrak{z}$-线性版本，定义见 7.9.7(i)。由 7.9.8，它作用在
$D_{HC\mathfrak{z}}$ 上。由于有显然的幺半函子
$\mathcal{H}\to\mathcal{H}_{\mathfrak{z}}$（见 7.9.7 的注），
$\mathcal{H}_{\mathfrak{z}}$ 包含 $\mathcal{P}$，故 $D_{HC\mathfrak{z}}$ 是
$\mathcal{P}$-模。如同 5.4.1，我们借助 Satake 等价把 $\mathcal{P}$ 换成
$\operatorname{Rep}{}^{L}G$，并把 $\operatorname{Rep}{}^{L}G$ 在
$D_{HC\mathfrak{z}}$ 上的相应作用记作 $\circledast$。另一方面，
$\mathcal{H}_{\mathfrak{z}}$ 的中心包含张量范畴
$\mathcal{M}(\operatorname{Aut}_ZO)^{fl}_{\mathfrak{z}}$，它由平坦
$\mathfrak{z}$-模配上 $\operatorname{Aut}_ZO$-作用构成（见 7.9.7(i)）。
$\mathcal{M}(\operatorname{Aut}_ZO)^{fl}_{\mathfrak{z}}$ 在 $D_{HC\mathfrak{z}}$ 上的
相应作用是显然的：对 $W\in\mathcal{M}(\operatorname{Aut}_ZO)^{fl}_{\mathfrak{z}}$、
$V\in D_{HC\mathfrak{z}}$，有
$W\circledast V=W\otimes V:=W\otimes_{\mathfrak{z}}V$。因此
$D_{HC\mathfrak{z}}$ 是
$(\operatorname{Rep}{}^{L}G,\mathcal{M}(\operatorname{Aut}_ZO)^{fl}_{\mathfrak{z}})$-双模。

设 $\mathfrak{F}$ 是 $\operatorname{Spec}\mathfrak{z}$ 上
$\operatorname{Aut}_ZO$-等变的 ${}^{L}G$-torsor。它给出张量函子
$\operatorname{Rep}{}^{L}G\to\mathcal{M}(\operatorname{Aut}_ZO)^{fl}_{\mathfrak{z}}$，
$V\mapsto V_{\mathfrak{F}}$，从而给出 $\operatorname{Rep}{}^{L}G$ 在
$D_{HC\mathfrak{z}}$ 上的相应作用。

\bdp{5.4.10} 我们在当前的 Harish-Chandra 设定中重复 5.4.2 的定义。即，
\emph{Hecke $\mathfrak{F}$-特征模}是一个 Harish-Chandra 模
$M\in\mathcal{M}_{HC\mathfrak{z}}$，连同自然同构
$\alpha_V:V\circledast M\xrightarrow{\sim}M\otimes V_{\mathfrak{F}}$
（$V\in\operatorname{Rep}{}^{L}G$），使得对任意
$V_1,V_2\in\operatorname{Rep}{}^{L}G$ 有
$\alpha_{V_1\otimes V_2}=\alpha_{V_1}\circ(V_1\circledast\alpha_{V_2})$。

现在可以表述 5.2.16 的精确版本。如同 5.2.15，我们的
$\mathcal{L}\in Z\operatorname{tors}_{\theta}(O)$（见 5.4.7）定义了局部
${}^{L}\mathfrak{g}$-oper 的模概形上的 $\operatorname{Aut}_ZO$-等变
${}^{L}G$-torsor\footnote{$\operatorname{Aut}_ZO$ 的作用来自等同
$\operatorname{Aut}_ZO=\operatorname{Aut}(O,\mathcal{L})$；见 4.6.6。}。
借助 Feigin--Frenkel 同构（80）把该概形与 $\operatorname{Spec}\mathfrak{z}$ 等同，
就得到 $\operatorname{Spec}\mathfrak{z}$ 上相应的
$\operatorname{Aut}_ZO$-等变 torsor $\mathfrak{F}_{\mathcal{L}}$。

\SourcePage{222}

从现在起，我们把 $\mathit{Vac}'$ 看作 $\mathcal{M}_{HC\mathfrak{z}}$ 的对象
（相对于固定真空向量的 $\operatorname{Aut}_ZO$-作用）。

\bdp{5.4.11} \emph{定理.} $\mathit{Vac}'$ 具有 Hecke
$\mathfrak{F}_{\mathcal{L}}$-特征模的典范结构。

这一定理蕴含 5.2.16。事实上，同构（282）是
$\operatorname{Aut}_ZO$-等变的，因为 $\operatorname{Aut}_ZO$ 通过结构的传输
作用在（282）的两端。

它将在何处证明？？？
% [Tutor Note: 原书 p.222 此处作「Where will it be proved???」，三个问号系底本
% 如此（作者未注明证明所在），译文照录。]

现在可以转向本节的主要结果。

\bdp{5.4.12} \emph{定理.} 定理 5.4.11 蕴含 5.4.5。

\emph{证明.} 我们将说明，一个适当的「局部化函子」${}^{L}\Delta$ 把局部图景
变为整体图景\footnote{读者可自行判断，这种疯狂之中是否有章法。}。

我们需要对 5.4.1 的设定稍作修改，以便使用 7.9 的「等变 Hecke 模式」。回忆在
整体定理 5.4.5 的表述中，我们固定了
$\mathcal{L}^{\mathrm{glob}}\in Z\operatorname{tors}_{\theta}(X)$（见 5.2.8），而在局部定理
5.4.11 中用的是 $\mathcal{L}^{\mathrm{loc}}\in Z\operatorname{tors}_{\theta}(O)$。考虑
4.4.15 中与 $\mathcal{L}^{\mathrm{glob}}$ 和 $\mathcal{L}^{\mathrm{loc}}$ 相应的概形
$X^{\wedge}_{Z}$ 与 $M^{\wedge}_{Z}$（它们是 5.4.1 所用概形 $X^{\wedge}$ 与
$M^{\wedge}$ 的平展 $Z$-覆叠）。回忆 $\operatorname{Aut}_ZO$ 作用在
$X^{\wedge}_{Z}$ 上，$\operatorname{Aut}_ZO\ltimes G(K)$ 作用在
$M^{\wedge}_{Z}$ 上（见 4.4.15）。有
$\operatorname{Aut}^{0}_{Z}O\backslash X^{\wedge}_{Z}=X$，且商叠
$(\operatorname{Aut}^{0}_{Z}O\ltimes G(O))\backslash M^{\wedge}_{Z}$ 等于
$\operatorname{Bun}_G\times X$。显然，在 5.4.1 中构造
$D(\operatorname{Bun}_G\times X)$ 上的 Hecke 作用时，可以把
$(M^{\wedge},\operatorname{Aut}O\ltimes G(K))$ 换成
$(M^{\wedge}_{Z},\operatorname{Aut}_ZO\ltimes G(K))$。

如同 5.1.1，设 $\lambda_{\mathcal{L}^{\mathrm{glob}}}$ 是 $\operatorname{Bun}_G$ 上对应于
$\mathcal{L}^{\mathrm{glob}}$ 的 Pfaffian 线丛。记
$\hat{\lambda}=\hat{\lambda}_{\mathcal{L}^{\mathrm{glob}}}$ 为它到 $M^{\wedge}_{Z}$ 的拉回。
$\operatorname{Aut}_ZO\ltimes G(K)$ 在 $M^{\wedge}_{Z}$ 上的作用以典范方式提升为
中心扩张 $\operatorname{Aut}_ZO\ltimes\widetilde{G(K)}$ 在 $\hat{\lambda}$ 上的作用
（见 4.4.16）。所以我们处于
\SourcePage{223}
7.9.6 的设定\footnote{很抱歉记号如此不一致：我们的 $M^{\wedge}_{Z}$、$X^{\wedge}$、
$\hat{\lambda}$、$\operatorname{Der}O$、$\operatorname{Aut}^{0}_{Z}O$、
$\widetilde{G(K)}$、$G(O)$ 分别是 7.9.6 的 $Y^{\wedge}$、$X^{\wedge}$、
$\mathcal{L}^{*}$、$\mathfrak{l}$、$P$、$G'$、$K$。}，于是有右 t-正合的
局部化函子
\[
{}^{L}\Delta:\ D_{HC}\to D(\operatorname{Bun}_G\times X)
\]

在 $\mathfrak{z}$-模的设定中也有相应的图景。事实上，按 7.9.7(ii)，考虑
$\mathcal{D}_X$-代数 $\mathfrak{z}_X$\footnote{任意 $\operatorname{Aut}O$-模 $V$
给出 $\mathcal{D}_X$-模 $V_X$，见 2.6.6。}（我们在 2.7 中已经用过它）以及相应的
范畴 $D(\operatorname{Bun}_G\times X,\mathfrak{z}_X)$，后者是
$\operatorname{Bun}_G\times X$ 上带 $\mathfrak{z}_X$-作用的 $\mathcal{D}$-模的
导出范畴（见 7.3.13）。它带有 $\mathcal{H}_{\mathfrak{z}}$ 的典范作用。有典范的
局部化函子
\[
{}^{L}\Delta_{\mathfrak{z}}:\ D_{HC\mathfrak{z}}\to
D(\operatorname{Bun}_G\times X,\mathfrak{z}_X)
\]
它是 $\mathcal{H}_{\mathfrak{z}}$-模的态射。上述 ${}^{L}\Delta$ 彼此相容（它们
与忘掉 $\mathfrak{z}$-作用可交换）。
% [Tutor Note: 原书 p.223 两处印作「Morhism」，系「Morphism」之误，译文按其本意
% 译为「态射」。]

现在我们的定理是下述事实的直接推论：

\begin{enumerate}[label=(\alph*),leftmargin=2.6em]
\item 有自然的等同
\begin{equation}\tag{283}
{}^{L}\Delta(\mathit{Vac}')=\Delta(\mathit{Vac}')
=M_{\mathcal{L}^{\mathrm{glob}}}\otimes\mathcal{O}_X
\end{equation}
使得 $\Delta(\mathit{Vac}')=\Delta_{\mathfrak{z}}(\mathit{Vac}')$ 上的
$\mathfrak{z}_X$-作用，与 $\mathfrak{z}_X$ 经由极大常值商
$\mathfrak{z}(X)\otimes\mathcal{O}_X=A_{{}^{L}\mathfrak{g}}(X)\otimes\mathcal{O}_X$
以及 $M_{\mathcal{L}^{\mathrm{glob}}}$ 上标准的 $A_{{}^{L}\mathfrak{g}}(X)$-模结构在
$M_{\mathcal{L}^{\mathrm{glob}}}\otimes\mathcal{O}_X$ 上的作用一致。证明见 7.14.9
（注意 $\mathfrak{z}_X$ 按结构的传输作用）。
\item 函子 ${}^{L}\Delta_{\mathfrak{z}}$ 是
$(\operatorname{Rep}{}^{L}G,\mathcal{M}(\operatorname{Aut}_ZO)^{fl}_{\mathfrak{z}})$-双模
的态射。事实上，它是 $\mathcal{H}_{\mathfrak{z}}$-模的态射。
\item 对任意 $W\in\mathcal{M}(\operatorname{Aut}_ZO)^{fl}_{\mathfrak{z}}$、
$T\in D(\operatorname{Bun}_G\times X,\mathfrak{z}_X)$，有
$W\circledast T=W_X\otimes_{\mathfrak{z}_X}T$，其中 $W_X$ 是对应于 $W$ 的
$\mathfrak{z}_X$-模。证明见 7.9.7(i)。
\end{enumerate}

\SourcePage{224}

\begin{enumerate}[label=(\alph*),leftmargin=2.6em]
\item[(d)] 对任意 $V\in\operatorname{Rep}{}^{L}G$，有典范等同
\[
(V_{\mathfrak{F}_{\mathcal{L}^{\mathrm{loc}}}})_X\otimes_{\mathfrak{z}_X}
(\mathfrak{z}(X)\otimes\mathcal{O}_X)\xrightarrow{\sim}
V_{\mathfrak{F}_{\mathcal{L}^{\mathrm{glob}}}}
\]
它与 $V$ 的张量积相容（这里 $\mathfrak{F}_{\mathcal{L}^{\mathrm{loc}}}$ 就是 5.4.10 的
$\mathfrak{F}_{\mathcal{L}}$）。\hfill$\Box$
\end{enumerate}

\bdsec{5.5. oper 的诞生}

本节假设定理 5.4.8 成立。我们首先说明，该定理蕴含 $\mathit{Vac}'$ 是 Hecke
$\mathcal{F}$-特征模，其中 $\mathcal{F}$ 是 $\operatorname{Spec}\mathfrak{z}$ 上
某个 $\operatorname{Aut}_ZO$-等变的 ${}^{L}G$-torsor。本节要点在于：
$\mathcal{F}$ 自然地来自局部 oper 的某个 $\operatorname{Aut}_ZO$-等变的
$\mathfrak{z}$-族。稍后我们将看到，
从 $\operatorname{Spec}\mathfrak{z}$ 到局部 oper 模空间的相应映射与
Feigin--Frenkel 同构一致，由此得到局部主定理。

\bdp{5.5.1} 对任意 $V\in\operatorname{Rep}{}^{L}G$，置
\begin{equation}\tag{284}
F_{\mathcal{H}}(V):=\operatorname{Hom}_{\widetilde{\mathfrak{g}\otimes K}}
(\mathit{Vac}',V\circledast\mathit{Vac}')=(V\circledast\mathit{Vac}')^{G(O)}.
\end{equation}

这是 $\operatorname{Aut}_ZO$-等变的 $\mathfrak{z}$-模\footnote{$V\circledast
\mathit{Vac}'$ 上的两个 $\mathfrak{z}$-模结构相同，因为 Hecke 函子是
$\mathfrak{z}$-线性的。}。由 5.4.8，它是自由 $\mathfrak{z}$-模，故
$F_{\mathcal{H}}(V)\in\mathcal{M}(\operatorname{Aut}_ZO)^{fl}_{\mathfrak{z}}$。
有典范同构
\begin{equation}\tag{285}
V\circledast\mathit{Vac}'=\mathit{Vac}'\otimes F_{\mathcal{H}}(V).
\end{equation}

由于 $\mathcal{M}(\operatorname{Aut}_ZO)^{fl}_{\mathfrak{z}}$ 的作用与 Hecke 作用
交换，我们得到典范等同
$F_{\mathcal{H}}(V_1\otimes V_2)=F_{\mathcal{H}}(V_1)\otimes F_{\mathcal{H}}(V_2)$，
这意味着
\begin{equation}\tag{286}
F_{\mathcal{H}}:\ \operatorname{Rep}{}^{L}G\to
\mathcal{M}(\operatorname{Aut}_ZO)^{fl}_{\mathfrak{z}}
\end{equation}
是幺半函子。

\bdp{5.5.2} \emph{引理.} 对任意 $V\in\operatorname{Rep}{}^{L}G$，自由
$\mathfrak{z}$-模 $F_{\mathcal{H}}(V)$ 有有限秩。

\emph{证明.} 由于 $F_{\mathcal{H}}$ 是幺半函子，$F_{\mathcal{H}}(V^{*})$ 在
幺半范畴的意义下对偶于 $F_{\mathcal{H}}(V)$（见 [Del91] 的 2.1.2）。若自由
$\mathfrak{z}$-模在幺半范畴的意义下有对偶，则其秩有限。\hfill$\Box$

\SourcePage{225}

设
\begin{equation}\tag{287}
F_{\mathfrak{F}_{\mathcal{L}}}:\ \operatorname{Rep}{}^{L}G\to
\mathcal{M}(\operatorname{Aut}_ZO)^{fl}_{\mathfrak{z}}
\end{equation}
为张量函子
$F_{\mathfrak{F}_{\mathcal{L}}}(V)=V_{\mathfrak{F}_{\mathcal{L}}}$（见 5.4.10）。

现在，我们的局部主定理 5.4.11 可重新表述如下。

\bdp{5.5.3} \emph{定理.} 幺半函子 $F_{\mathcal{H}}$ 与
$F_{\mathfrak{F}_{\mathcal{L}}}$ 典范同构。

我们将说明，$F_{\mathcal{H}}$ 确实来自某个由 $\operatorname{Spec}\mathfrak{z}$
参数化的、典范定义的局部 oper 族。首先核对 $F_{\mathcal{H}}$ 确实来自
$\operatorname{Spec}\mathfrak{z}$ 上的一个 ${}^{L}G$-torsor。

\bdp{5.5.4} \emph{命题.} 幺半函子 $F_{\mathcal{H}}$ 是张量函子，即它与交换
约束相容。

证明分两步。首先把相容同构
$F_{\mathcal{H}}(V_1)\otimes F_{\mathcal{H}}(V_2)\simeq
F_{\mathcal{H}}(V_1\otimes V_2)$ 写成（扭）$\mathcal{D}$-模截面的卷积积
（见 5.5.5、5.5.6）。然后利用卷积的融合图景，说明它交换（见 ???）。
% [Tutor Note: 原书 p.225 此处交叉引用印作「???」，底本未补，译文照录。]

\bdp{5.5.5} 我们借助 Satake 等价 $h$（见（273）），把 ${}^{L}G$-模的张量范畴
换成仿射 Grassmann 流形上 $\mathcal{D}$-模的张量范畴。对
$P\in\mathcal{P}^{\natural}$，置 $F_{\mathcal{H}}(P):=F_{\mathcal{H}}(hP)$。
于是（见（282））
\begin{equation}\tag{288}
F_{\mathcal{H}}(P)=\Gamma(\mathcal{GR},P\lambda_{\mathcal{L}}^{-1})^{G(O)}.
\end{equation}

\emph{注.} 回忆 $P$ 是「超」$\mathcal{D}$-模，$\lambda_{\mathcal{L}}$ 是「超」
线丛。然而它们的奇偶性相同（都等于 $\mathcal{GR}$ 各连通分支的奇偶性），故
$P\lambda_{\mathcal{L}}^{-1}$ 是通常的偶层。这些「超」的微妙之处在我们转向
交换约束时才会起作用。

为描述相容同构
$F_{\mathcal{H}}(P_1)\otimes F_{\mathcal{H}}(P_2)\simeq
F_{\mathcal{H}}(P_1\circledast P_2)$，考虑 $\mathcal{O}^{!}$-模的积分态射
（我们用 5.3.5 的记号；关于积分见 7.11.16 的（??））
\begin{equation}\tag{289}
i_m:\ m\cdot(P_1\boxtimes'P_2)\to P_1\circledast P_2.
\end{equation}
% [Tutor Note: 原书 p.225 印作「7.11.16 (??)」，括号内编号未补全，译文照录。]

\SourcePage{226}

$\mathcal{GR}$ 上的线丛 $\lambda_{\mathcal{L}}$ 是 $G(O)$-等变的，它经由
$m:G(K)\times_{G(O)}\mathcal{GR}\to\mathcal{GR}$ 的拉回被典范地等同于「扭积」
$\lambda_{\mathcal{L}}\boxtimes'\lambda_{\mathcal{L}}$\footnote{这是因为按定义，
$\lambda_{\mathcal{L}}$ 来自 $G(K)$ 的一个中心扩张，该扩张在 $G(O)$ 上有分裂。}。
于是用 $\lambda_{\mathcal{L}}$ 扭 $i_m$，得到态射
\[
m\cdot((P_1\lambda_{\mathcal{L}}^{-1})\boxtimes'(P_2\lambda_{\mathcal{L}}^{-1}))
\to(P_1\circledast P_2)\lambda_{\mathcal{L}}^{-1}.
\]

过渡到 $G(O)$-不变截面，就得到\emph{卷积映射}（注意 $G(O)$-不变性允许我们
忽略扭）
\begin{equation}\tag{290}
*:\ \Gamma(\mathcal{GR},P_1\lambda_{\mathcal{L}}^{-1})^{G(O)}\otimes
\Gamma(\mathcal{GR},P_2\lambda_{\mathcal{L}}^{-1})^{G(O)}\to
\Gamma(\mathcal{GR},(P_1\circledast P_2)\lambda_{\mathcal{L}}^{-1})^{G(O)}.
\end{equation}

\bdp{5.5.6} \emph{引理.} 卷积映射与相容同构
$F_{\mathcal{H}}(P_1)\otimes F_{\mathcal{H}}(P_2)\simeq
F_{\mathcal{H}}(P_1\circledast P_2)$ 一致。

\emph{证明.} 考虑典范同构（作用约束）
$a:P_1\circledast(P_2\circledast\mathit{Vac}')\xrightarrow{\sim}
(P_1\circledast P_2)\circledast\mathit{Vac}'$。对
$f\in\operatorname{Hom}(\mathit{Vac}',P_1\circledast\mathit{Vac}')$、
$g\in\operatorname{Hom}(\mathit{Vac}',P_2\circledast\mathit{Vac}')$，相容同构
把 $f\otimes g$ 送到 $(P_1\circledast g)\circ f$。\hfill$\Box$

\bdsec{5.6. 重整化泛包络代数}

\bdp{5.6.1} 设 $A$ 为 $\widetilde{\mathfrak{g}\otimes K}$ 的完备化泛包络代数。
由 3.6.2，$A$ 是 $\mathbb{C}[h]$ 上的平坦代数，$h:=\mathbf{1}-1$，且
$A/hA=\overline{U}'$。

$A$ 上的自然拓扑在
$A[h^{-1}]:=A\otimes_{\mathbb{C}[h]}\mathbb{C}[h,h^{-1}]$ 上诱导一个拓扑；
事实上这是归纳极限拓扑（把 $A[h^{-1}]$ 表示成 $A\to A\to\ldots$ 的归纳极限，
其中每个箭头都是乘以 $h$）。

设 $I\subset\mathfrak{Z}$ 是 3.6.5 中的理想。记 $J$ 为
$I\overline{U}'\subset\overline{U}'=A/hA$ 在 $A$ 中的原像（$I\overline{U}'$ 按
拓扑意义理解，即 $I\overline{U}'$ 是由 $I$ 生成的 $\overline{U}'$ 的闭理想）。
$J$ 是 $A$ 中包含 $hA$ 的闭理想。记 $A^{\natural}$ 为递增序列
$A\subset h^{-1}J\subset h^{-2}J^{2}\subset\ldots$ 的并，其中 $J^{k}$ 按拓扑
意义理解。最后置 $U^{\natural}:=A^{\natural}/hA^{\natural}$。

\SourcePage{227}

$A^{\natural}$ 是 $\mathbb{C}[h]$ 上的拓扑代数（拓扑由 $A[h^{-1}]$ 诱导）。故
$U^{\natural}$ 是拓扑 $\mathbb{C}$-代数（$U^{\natural}$ 带商拓扑）。

\bdp{5.6.2} 置 $\mathit{Vac}_A=A/A(\mathfrak{g}\otimes O)$，其中
$A(\mathfrak{g}\otimes O)$ 表示 $A$ 中由 $\mathfrak{g}\otimes O$ 生成的闭左理想。
$I$ 在 $\mathit{Vac}'=\mathit{Vac}_A/h\mathit{Vac}_A$ 上平凡地作用。由于
$\mathit{Vac}_A$ 是平坦 $\mathbb{C}[h]$-模，$A^{\natural}$ 作用在
$\mathit{Vac}_A$ 上。因此 $U^{\natural}$ 作用在 $\mathit{Vac}'$ 上。

\bdp{5.6.3} 记 $U^{\natural}_0$ 为 $A$ 在 $U^{\natural}$ 中的像。
$U^{\natural}_0$ 是 $U^{\natural}$ 的子代数。我们给 $U^{\natural}_0$ 配上诱导
拓扑。映射 $A\to U^{\natural}_0$ 经由 $A/hA=\overline{U}'$ 分解，事实上经由
$\overline{U}'/I\overline{U}'$ 分解。于是得到满的连续同态
$f:\overline{U}'/I\overline{U}'\to U^{\natural}_0$。$f$ 大概是同胚。

无论如何，$f$ 诱导拓扑同构
$\mathfrak{z}=\mathfrak{Z}/I\xrightarrow{\sim}f(\mathfrak{z})$（用
$U^{\natural}$ 在 $\mathit{Vac}'$ 上的作用）。我们将把 $\mathfrak{z}$ 与
$f(\mathfrak{z})$ 等同。

\bdp{5.6.4} 设 $J_I\subset A$ 表示 $I\subset\overline{U}'=A/hA$ 的原像。记
$U^{\natural}_1$ 为 $h^{-1}J_I$ 在 $U^{\natural}$ 中的像。给 $U^{\natural}_1$ 配上
由 $U^{\natural}$ 诱导的拓扑。拓扑代数 $U^{\natural}$ 由 $U^{\natural}_1$ 生成。

\bdp{5.6.5} \emph{引理.}
\begin{enumerate}[label=(\roman*),leftmargin=2.2em]
\item $U^{\natural}_1$ 是 $U^{\natural}$ 的李子代数。
\item $U^{\natural}_0$ 是 $U^{\natural}_1$ 的理想。
\item $\mathfrak{z}U^{\natural}_1\subset U^{\natural}_1$，
$U^{\natural}_1\mathfrak{z}\subset U^{\natural}_1$。
\item $[U^{\natural}_1,\mathfrak{z}]\subset\mathfrak{z}$。
\end{enumerate}

\emph{证明.} 我们要用到 3.6.2 中定义的 Hayashi 括号 $\{\ ,\ \}$ 的若干性质。
(i) 由包含 $[J_I,J_I]\subset hJ_I$ 推出，后者显然，因为 $\{I,I\}\subset I$
（见 3.6.4 (i)）。(ii) 与 (iii) 显然。(iv) 显然，因为
$\{I,\mathfrak{Z}\}\subset\{\mathfrak{Z},\mathfrak{Z}\}\subset\mathfrak{Z}$。
\hfill$\Box$

\bdp{5.6.6} 由 5.6.5 可知，$U^{\natural}_1/U^{\natural}_0$ 是 $\mathfrak{z}$ 上的
拓扑李代数胚。乘以 $h^{-1}$ 定义了映射 $J_I\to A^{\natural}$，它诱导李代数胚
态射
\begin{equation}\tag{291}
I/I^{2}=J_I/(J_I^{2}+hA)\to U^{\natural}_1/U^{\natural}_0
\end{equation}

\SourcePage{228}

（$I/I^{2}$ 上代数胚结构的定义见 3.6.5。）态射（291）连续且满。事实上它是
拓扑同构（见 ???）。
% [Tutor Note: 原书 p.228 此处交叉引用印作「???」，底本未补，译文照录。]

\bdp{5.6.7} 记 $U^{\flat}_i$ 为 $U^{\natural}_i$ 中消灭 $\mathit{Vac}'$ 真空向量的
元素之集，$i=0,1$。引理 5.6.5 在把 $U^{\natural}_i$ 换成 $U^{\flat}_i$
（$i=0,1$）后仍然成立。故 $U^{\flat}_1/U^{\flat}_0$ 是 $\mathfrak{z}$ 上的拓扑
李代数胚。自然映射 $U^{\flat}_1/U^{\flat}_0\to U^{\natural}_1/U^{\natural}_0$ 是
拓扑同构。于是（291）诱导满的连续李代数胚态射
\begin{equation}\tag{292}
I/I^{2}\to U^{\flat}_1/U^{\flat}_0.
\end{equation}

\bdp{5.6.8} 设 $V$ 是拓扑 $U^{\natural}$-模（在我们关心的应用中 $V$ 是离散的）。
则 $V^{\mathfrak{g}\otimes O}$ 是李代数胚 $I/I^{2}$ 上的（左）拓扑模。事实上，
首先 $V^{\mathfrak{g}\otimes O}$ 是 $\mathfrak{z}$-模。其次，
$V^{\mathfrak{g}\otimes O}=\{v\in V\mid U^{\flat}_0v=0\}$，故李代数
$U^{\flat}_1/U^{\flat}_0$ 作用在 $V^{\mathfrak{g}\otimes O}$ 上。若
$v\in V^{\mathfrak{g}\otimes O}$、$z\in\mathfrak{z}$、
$a\in U^{\flat}_1/U^{\flat}_0$，则 $a(zv)-z(av)=\partial_a(z)v$，其中
$\partial_a\in\operatorname{Der}\mathfrak{z}$ 按 $U^{\flat}_1/U^{\flat}_0$ 上的
代数胚结构对应于 $a$。于是 $V^{\mathfrak{g}\otimes O}$ 是代数胚
$U^{\flat}_1/U^{\flat}_0$ 上的模。利用（292）可见 $V^{\mathfrak{g}\otimes O}$ 是
李代数胚 $I/I^{2}$ 上的模。

\bdp{5.6.9} 由（89），有连续李代数态射
$\operatorname{Der}O\to h^{-1}J_I\subset A[h^{-1}]$，使得
$L_n\mapsto h^{-1}\widetilde{\mathcal{L}}_n$，$n\ge -1$。它诱导连续李代数态射
$\operatorname{Der}O\to U^{\flat}_1\subset U^{\natural}$。另一方面，我们在
3.6.16 中定义了典范态射 $\operatorname{Der}O\to I/I^{2}$。显然下图交换：
\[
\begin{array}{ccc}
\operatorname{Der}O & \longrightarrow & U^{\flat}_1\\[2pt]
\downarrow & & \downarrow\\[2pt]
I/I^{2} & \longrightarrow & U^{\flat}_1/U^{\flat}_0
\end{array}
\]

\emph{注.} 态射 $\operatorname{Der}O\to U^{\flat}_1/U^{\flat}_0$ 把
$\operatorname{Der}O$ 同胚地映到其像上。由于 $U^{\flat}_1/U^{\flat}_0$ 连续
作用在 $\mathfrak{z}\subset U^{\natural}_0$ 上，这可由
\SourcePage{229}
关于态射 $\operatorname{Der}O\to\operatorname{Der}\mathfrak{z}$ 的相应论断推出，
而后者显然（看 $\mathfrak{z}$ 的 Sugawara 元素）。

\bdp{5.6.10} 设我们处于 5.6.8 的情形。由 5.6.9，$\operatorname{Der}O$ 经由
态射 $\operatorname{Der}O\to U^{\natural}$ 作用在 $V$ 上，子空间
$V^{\mathfrak{g}\otimes O}$ 是 $\operatorname{Der}O$-不变的，且
$\operatorname{Der}O$ 在 $V^{\mathfrak{g}\otimes O}$ 上的作用与来自态射
$\operatorname{Der}O\to I/I^{2}$ 的作用一致。

\bdp{5.6.11} \emph{注.} 2.5.1 中 $\widetilde{\mathfrak{g}\otimes K}$ 的定义涉及
由（18）定义的「临界」数量积 $c$。设想考虑对应于 $\lambda c$ 的中心扩张
$0\to\mathbb{C}\to(\widetilde{\mathfrak{g}\otimes K})_{\lambda}\to
\mathfrak{g}\otimes K\to 0$，$\lambda\in\mathbb{C}^{*}$。记 $A_{\lambda}$ 为
$(\widetilde{\mathfrak{g}\otimes K})_{\lambda}$ 的完备化泛包络代数。若把 $A$ 与
$h=\mathbf{1}-1$ 换成 $A_{\lambda}$ 与
$h_{\lambda}:=\mathbf{1}_{\lambda}-\lambda^{-1}$，则 $U^{\natural}$ 的构造与映射
（291）仍然成立；这里 $\mathbf{1}_{\lambda}$ 表示
$1\in\mathbb{C}\subset(\widetilde{\mathfrak{g}\otimes K})_{\lambda}$。记
$U^{\natural}_{\lambda}$ 与 $f_{\lambda}$ 为对应于 $\lambda$ 的
$U^{\natural}$ 与（291）的类似物。可以利用典范同构
$\widetilde{\mathfrak{g}\otimes K}\xrightarrow{\sim}
(\widetilde{\mathfrak{g}\otimes K})_{\lambda}$（它使
$\mathbf{1}\mapsto\lambda\cdot\mathbf{1}_{\lambda}$）把 $A_{\lambda}$ 与
$U^{\natural}_{\lambda}$ 分别与 $A$ 与 $U^{\natural}$ 等同。此时 $f_{\lambda}$
\emph{确实}依赖于 $\lambda$：事实上 $f_{\lambda}=\lambda f_1$。

% 尾部接 p.230（本页末句完结于 f_λ = λ f_1，5.6 节尚未结束）。
